\chapter{单自由度系统动力学建模}

\intro{
	最简单的系统往往蕴含着最深刻的物理。
	
	\rightline{——Enrico Fermi}
}

\section{弹簧-质量-阻尼系统}

单自由度（SDOF）系统是最基础的力学系统，其典型实例为弹簧-质量-阻尼系统。

\subsection{运动方程}

由 Newton 第二定律，位移 $x(t)$ 满足：
\[
m \ddot{x}(t) + c \dot{x}(t) + k x(t) = F(t)
\]
其中各项物理意义：
\begin{itemize}
	\item $m \ddot{x}$——惯性力（质量 $m$ 乘加速度）；
	\item $c \dot{x}$——阻尼力（线性黏性阻尼，$c$ 为阻尼系数）；
	\item $k x$——弹性恢复力（Hooke 定律，$k$ 为刚度）；
	\item $F(t)$——外激励力。
\end{itemize}

标准化为：
\begin{myformula}
	\ddot{x} + 2 \zeta \omega_n \dot{x} + \omega_n^2 x = \frac{F(t)}{m}
\end{myformula}
其中：
\begin{myformula}
	\omega_n = \sqrt{\frac{k}{m}} \quad \text{（固有频率）}, \qquad
	\zeta = \frac{c}{2\sqrt{mk}} \quad \text{（阻尼比）}
\end{myformula}

\section{自由振动}

\subsection{无阻尼自由振动（$\zeta = 0$）}

方程为 $\ddot{x} + \omega_n^2 x = 0$，其通解为：
\begin{myformula}
	x(t) = A \cos(\omega_n t) + B \sin(\omega_n t) = X \cos(\omega_n t - \varphi)
\end{myformula}
其中振幅 $X = \sqrt{A^2 + B^2} = \sqrt{x_0^2 + (v_0/\omega_n)^2}$，初相位 $\varphi = \arctan(v_0/(\omega_n x_0))$。

\subsection{欠阻尼自由振动（$0 < \zeta < 1$）}

特征方程 $s^2 + 2\zeta\omega_n s + \omega_n^2 = 0$ 的根为：
\[
s_{1,2} = -\zeta \omega_n \pm i \omega_d, \quad \omega_d = \omega_n \sqrt{1 - \zeta^2}
\]
其中 $\omega_d$ 为\textbf{有阻尼固有频率}。

通解为：
\begin{myformula}
	x(t) = e^{-\zeta \omega_n t} (A \cos \omega_d t + B \sin \omega_d t) = X e^{-\zeta \omega_n t} \cos(\omega_d t - \varphi)
\end{myformula}

\begin{remark}
	欠阻尼振动的振幅按指数规律衰减。对数衰减率：
	\[
	\delta = \ln \frac{x_k}{x_{k+1}} = \frac{2\pi \zeta}{\sqrt{1 - \zeta^2}}
	\]
	实验上，通过测量相邻周期的振幅比可确定阻尼比 $\zeta$。
\end{remark}

\subsection{临界阻尼（$\zeta = 1$）与过阻尼（$\zeta > 1$）}

\begin{itemize}
	\item \textbf{临界阻尼}：$s_{1,2} = -\omega_n$（重根），解为：
	\[
	x(t) = (A + Bt) e^{-\omega_n t}
	\]
	系统以最快速度回到平衡位置而不发生振荡。
	\item \textbf{过阻尼}：$s_{1,2} = -\zeta \omega_n \pm \omega_n \sqrt{\zeta^2 - 1}$（两相异负实根），解为指数衰减之和，不产生振动。
\end{itemize}

\section{受迫振动}

\subsection{简谐激励下的稳态响应}

设激励力 $F(t) = F_0 \cos(\omega t)$（或复数形式 $F_0 e^{i\omega t}$），稳态解具有形式：
\begin{myformula}
	x_p(t) = X \cos(\omega t - \theta)
\end{myformula}
代入运动方程得幅值和相位：
\begin{myformula}
	X = \frac{F_0 / k}{\sqrt{(1 - r^2)^2 + (2\zeta r)^2}}, \qquad
	\theta = \arctan \frac{2\zeta r}{1 - r^2}
\end{myformula}
其中 $r = \omega/\omega_n$ 为\textbf{频率比}。

\begin{mydefinition}{放大因子与品质因子}
	放大因子 $M = X / (F_0/k)$ 表示动态位移与静位移之比。当 $\zeta$ 较小时：
	\[
	M_{\max} \approx \frac{1}{2\zeta} = Q
	\]
	$Q$ 称为\textbf{品质因子}，反映共振尖锐程度。
\end{mydefinition}

\begin{theorem}{共振条件}
	幅值极大值发生在：
	\[
	r_{\text{peak}} = \sqrt{1 - 2\zeta^2}
	\]
	当 $\zeta \ll 1$ 时，$r_{\text{peak}} \approx 1$，即激励频率接近固有频率时发生共振。
\end{theorem}

\begin{proof}
	令 $\frac{d}{dr} [(1 - r^2)^2 + (2\zeta r)^2]^{-1/2} = 0$：
	\[
	-\frac{1}{2}[(1 - r^2)^2 + (2\zeta r)^2]^{-3/2} \cdot [-4r(1 - r^2) + 8\zeta^2 r] = 0
	\]
	即 $-4r(1 - r^2) + 8\zeta^2 r = 0$，整理得 $r(1 - r^2 - 2\zeta^2) = 0$，故 $r_{\text{peak}} = \sqrt{1 - 2\zeta^2}$。
\end{proof}

\subsection{复频响应函数}

使用复数表示：$F(t) = F_0 e^{i\omega t}$，$x_p = \hat{X} e^{i\omega t}$，代入得：
\begin{myformula}
	(-\omega^2 m + i \omega c + k) \hat{X} = F_0
\end{myformula}
定义\textbf{机械阻抗} $Z(\omega) = -\omega^2 m + i\omega c + k$，则：
\begin{myformula}
	H(\omega) = \frac{\hat{X}}{F_0} = \frac{1}{Z(\omega)} = \frac{1}{k - \omega^2 m + i \omega c}
\end{myformula}
$H(\omega)$ 称为\textbf{频率响应函数（FRF）}，在实验模态分析中具有核心地位。

\section{任意激励下的响应——脉冲响应与卷积}

\subsection{单位脉冲响应函数}

当 $F(t) = \delta(t)$（Dirac $\delta$ 函数）时，系统的响应 $h(t)$ 称为\textbf{单位脉冲响应函数}。对于欠阻尼 SDOF 系统：
\begin{myformula}
	h(t) = \frac{1}{m \omega_d} e^{-\zeta \omega_n t} \sin(\omega_d t), \quad t \ge 0
\end{myformula}

\subsection{Duhamel 积分（卷积积分）}

对任意激励 $F(t)$，响应可由脉冲响应的叠加求得：
\begin{myformula}
	x(t) = \int_{0}^{t} F(\tau) \, h(t - \tau) \, d\tau = \int_{0}^{t} h(\tau) \, F(t - \tau) \, d\tau
\end{myformula}

\begin{proof}
	将 $F(t)$ 视作一系列脉冲 $\delta$ 函数的叠加 $F(t) = \int_{0}^{\infty} F(\tau) \delta(t - \tau) d\tau$。由线性系统的叠加原理和时不变性，输出为各脉冲响应的叠加：
	\[
	x(t) = \int_{0}^{\infty} F(\tau) h(t - \tau) d\tau
	\]
	因 $h(t - \tau) = 0$（$t < \tau$），故积分上限为 $t$。
\end{proof}

\section{能量方法}

\subsection{动能与势能}

SDOF 系统的能量：
\begin{itemize}
	\item 动能：$T = \frac{1}{2} m \dot{x}^2$
	\item 势能（弹性）：$V = \frac{1}{2} k x^2$
	\item 耗散能率：$P_d = c \dot{x}^2$
\end{itemize}

能量平衡方程（对无激励的耗散系统积分运动方程）：
\begin{myformula}
	\frac{d}{dt}(T + V) = -c \dot{x}^2 \le 0
\end{myformula}
即机械能随时间单调递减。

\subsection{等效质量与等效刚度——能量法}

对于连续体或复杂结构，用能量等效原理确定等效参数。

\textbf{Rayleigh 能量法：} 假设振型函数 $\phi(x)$，位移场 $u(x,t) = \phi(x) q(t)$，则：
\begin{myformula}
	T = \frac{1}{2} m_{\text{eq}} \dot{q}^2, \qquad V = \frac{1}{2} k_{\text{eq}} q^2
\end{myformula}
其中：
\begin{myformula}
	m_{\text{eq}} = \int m(x) \phi^2(x) \, dx, \qquad
	k_{\text{eq}} = \int EI(x) \left[ \phi''(x) \right]^2 \, dx
\end{myformula}
固有频率估计：$\omega_n^2 = k_{\text{eq}} / m_{\text{eq}}$。

\section{状态空间实现}

将 SDOF 系统改写为一阶方程组——状态空间形式：
\begin{myformula}
	\begin{pmatrix} \dot{x}_1 \\ \dot{x}_2 \end{pmatrix}
	=
	\begin{pmatrix} 0 & 1 \\ -\omega_n^2 & -2\zeta\omega_n \end{pmatrix}
	\begin{pmatrix} x_1 \\ x_2 \end{pmatrix}
	+
	\begin{pmatrix} 0 \\ 1/m \end{pmatrix} F(t)
\end{myformula}
其中 $x_1 = x$，$x_2 = \dot{x}$。此形式可直接代入 MATLAB 的 \texttt{ode45} 等求解器：

\begin{example}{MATLAB 求解 SDOF 系统}
	\begin{verbatim}
	m = 1;  c = 0.1;  k = 10;
	odefun = @(t, x) [x(2); (-c*x(2) - k*x(1))/m];
	[t, X] = ode45(odefun, [0 20], [1; 0]);
	plot(t, X(:,1));  % 位移时程
	\end{verbatim}
\end{example}
